% ---------- Language and typography ----------
% NOTE: in babel the LAST language listed is the main one.  `main=english`
% keeps the document in English (Contents / Table / mod) while still allowing
% \selectlanguage{spanish} for occasional Spanish quotations.
\usepackage[main=english,spanish]{babel}
\usepackage[T1]{fontenc}
\usepackage[utf8]{inputenc}
\usepackage{lmodern}
\usepackage{microtype}

% ---------- Page layout ----------
\usepackage[a4paper,top=2.6cm,bottom=2.5cm,left=2.7cm,right=2.5cm]{geometry}
\usepackage{setspace}
\setstretch{1.12}
\setlength{\parindent}{0pt}
\setlength{\parskip}{6pt}

% ---------- Mathematics ----------
\usepackage{amsmath,amssymb}

% ---------- Tables ----------
\usepackage{booktabs}
\usepackage{tabularx}
\usepackage{array}
\usepackage{longtable}

% ---------- Figures and graphics ----------
\usepackage{graphicx}
\usepackage{float}
\usepackage{caption}
\usepackage{subcaption}
\usepackage{placeins}
\usepackage{tikz}

% ---------- Code ----------
\usepackage{xcolor}
\usepackage{listings}

% ---------- Quotations (required by biblatex when babel is loaded) ----------
\usepackage{csquotes}

% ---------- Links ----------
\usepackage[hidelinks]{hyperref}
\usepackage[nameinlink,noabbrev]{cleveref}

% ---------- Headers / footers ----------
\usepackage{fancyhdr}

% ---------- Lists ----------
\usepackage{enumitem}

% ---------- References ----------
\usepackage[backend=biber,style=ieee]{biblatex}
\addbibresource{references.bib}

% ============================================================
% DOCUMENT INFORMATION — CHANGE THESE VALUES FOR EACH REPORT
% ============================================================
% ============================================================
% DOCUMENT INFORMATION — CHANGE THESE VALUES FOR EACH REPORT
% ============================================================

\newcommand{\university}{Yachay Tech University}
\newcommand{\faculty}{School of Mathematical and Computational Sciences}

\newcommand{\reporttitle}{Secure Wireless Messaging for IoT Devices}
\newcommand{\reportsubtitle}{DHE-PSK key exchange and Encrypt-then-MAC on two ESP32 boards}

% Shown on the cover above the title
\newcommand{\reportkind}{Project 1 -- Technical Report}

% TODO: add the name of the second team member (Member B) next to the author, e.g. "Daniel Troya A. and Name Surname"
\newcommand{\reportauthor}{Daniel Troya, Jose Reinoso}
\newcommand{\reportcourse}{Cryptography}
\newcommand{\reportprofessor}{PhD. Christian Iza}

\newcommand{\reportlocation}{Urcuqu\'i, Ecuador}
\newcommand{\reportdate}{October 2026}


% PDF metadata (must come after labels.tex so the macros are defined)
\hypersetup{
    pdftitle={\reporttitle},
    pdfauthor={\reportauthor},
    pdfsubject={\reportcourse},
    pdfcreator={LaTeX}
}

% ---------- Visual identity ----------
\definecolor{accent}{HTML}{243447}
\definecolor{schoolred}{HTML}{D3392A} % Institutional red (taken from the School logo)
\definecolor{rulegray}{HTML}{9AA1A9}
\definecolor{codebackground}{HTML}{F5F5F5}
\definecolor{codecomment}{HTML}{666666}
\definecolor{codekeyword}{HTML}{333333}

% ============================================================
% PAGE STYLE
% ============================================================

\pagestyle{fancy}
\fancyhf{}

% Header marks.  The original template used \leftmark, which fancyhdr resolves
% to the LAST section started on the page: page 1 was headed "4 Methodology"
% even though it was almost entirely Introduction.  \markright + \rightmark
% resolves to the FIRST section on the page instead, which is the right name.
% \nouppercase cancels the forced all-caps of the default mark.
\renewcommand{\sectionmark}[1]{\markright{\thesection\quad #1}}
\renewcommand{\subsectionmark}[1]{} % keep subsections out of the header

\fancyhead[L]{\small\textsc{\reporttitle}}
\fancyhead[R]{\small\textsc{\nouppercase{\rightmark}}}

\fancyfoot[C]{\small\thepage}

\renewcommand{\headrulewidth}{0.4pt}
\renewcommand{\footrulewidth}{0pt}

\setlength{\headheight}{14pt}

% ============================================================
% SECTION STYLE
% ============================================================

\usepackage{titlesec}

\titleformat{\section}
  {\Large\bfseries}
  {\thesection}
  {0.7em}
  {}

\titleformat{\subsection}
  {\large\bfseries}
  {\thesubsection}
  {0.7em}
  {}

\titleformat{\subsubsection}
  {\normalsize\bfseries}
  {\thesubsubsection}
  {0.7em}
  {}

\titlespacing*{\section}{0pt}{24pt}{10pt}
\titlespacing*{\subsection}{0pt}{18pt}{8pt}
\titlespacing*{\subsubsection}{0pt}{12pt}{6pt}

% ============================================================
% CAPTIONS
% ============================================================

\captionsetup{
    font=small,
    labelfont=bf,
    labelsep=period,
    justification=centering
}

% ============================================================
% LISTINGS — CODE STYLE
% ============================================================

\lstdefinestyle{formalcode}{
    backgroundcolor=\color{codebackground},
    basicstyle=\ttfamily\small,
    keywordstyle=\bfseries\color{codekeyword},
    commentstyle=\itshape\color{codecomment},
    stringstyle=\color{black},
    numbers=left,
    numberstyle=\tiny\color{codecomment},
    numbersep=10pt,
    frame=single,
    framerule=0.4pt,
    rulecolor=\color{black!25},
    breaklines=true,
    breakatwhitespace=false,
    showstringspaces=false,
    tabsize=4,
    keepspaces=true,
    columns=fullflexible,
    % UTF-8 support: without this, pasting console output containing accents
    % or ñ (very common in lab transcripts) breaks the compilation.
    inputencoding=utf8,
    extendedchars=true,
    literate=%
        {á}{{\'a}}1 {é}{{\'e}}1 {í}{{\'i}}1 {ó}{{\'o}}1 {ú}{{\'u}}1
        {Á}{{\'A}}1 {É}{{\'E}}1 {Í}{{\'I}}1 {Ó}{{\'O}}1 {Ú}{{\'U}}1
        {ñ}{{\~n}}1 {Ñ}{{\~N}}1 {ü}{{\"u}}1 {Ü}{{\"U}}1
        {¿}{{?`}}1 {¡}{{!`}}1 {°}{{$^\circ$}}1 {€}{{\texteuro}}1
}

\lstset{style=formalcode}

% ============================================================
% CUSTOM COMMANDS
% ============================================================

% Highlighted technical term
\newcommand{\term}[1]{\textbf{#1}}

% Simple figure command:  \reportfigure[width]{file}{caption}{label}
% The label argument lets you reference the figure with \cref{fig:...}.
\newcommand{\reportfigure}[4][0.90\textwidth]{%
    \begin{figure}[H]
        \centering
        \includegraphics[width=#1]{#2}
        \caption{#3}
        \label{#4}
    \end{figure}%
}

% ============================================================
% CRYPTOGRAPHY TOOLKIT
% Theorem environments, algorithms, notation macros and protocol
% diagrams.  See cheatsheet.pdf for a worked example of everything.
% ============================================================
% ============================================================
% CRYPTOGRAPHY TOOLKIT
% Theorem environments, algorithm layout, notation macros and
% TikZ styles for two-party protocol diagrams.
% This file is loaded at the end of packages.tex.
% ============================================================

% ---------- Packages ----------
\usepackage{amsthm}
\usepackage{mathtools}
\usepackage{algorithm}
\usepackage{algpseudocode}

\usetikzlibrary{arrows.meta,positioning,calc}

% ============================================================
% THEOREM ENVIRONMENTS
% All share one counter, numbered within the section
% (Definition 3.1, Theorem 3.2, ...), which is the usual
% convention in cryptography texts.
% ============================================================

\theoremstyle{definition}
\newtheorem{definition}{Definition}[section]
\newtheorem{construction}[definition]{Construction}
\newtheorem{example}[definition]{Example}
\newtheorem{remark}[definition]{Remark}

\theoremstyle{plain}
\newtheorem{theorem}[definition]{Theorem}
\newtheorem{lemma}[definition]{Lemma}
\newtheorem{corollary}[definition]{Corollary}
\newtheorem{proposition}[definition]{Proposition}
\newtheorem{claim}[definition]{Claim}

% ---------- cleveref names ----------
\crefname{definition}{Definition}{Definitions}
\Crefname{definition}{Definition}{Definitions}
\crefname{theorem}{Theorem}{Theorems}
\Crefname{theorem}{Theorem}{Theorems}
\crefname{lemma}{Lemma}{Lemmas}
\Crefname{lemma}{Lemma}{Lemmas}
\crefname{corollary}{Corollary}{Corollaries}
\Crefname{corollary}{Corollary}{Corollaries}
\crefname{construction}{Construction}{Constructions}
\Crefname{construction}{Construction}{Constructions}
\crefname{algorithm}{Algorithm}{Algorithms}
\Crefname{algorithm}{Algorithm}{Algorithms}
\crefname{lstlisting}{Listing}{Listings}
\Crefname{lstlisting}{Listing}{Listings}

% ============================================================
% ALGORITHM LAYOUT
% ============================================================

\algrenewcommand\algorithmicrequire{\textbf{Input:}}
\algrenewcommand\algorithmicensure{\textbf{Output:}}
\algrenewcommand\algorithmicforall{\textbf{for each}}

% Vertical guide lines for nested blocks
\algrenewcommand\algorithmicindent{1.1em}
\makeatletter
\renewcommand{\ALG@name}{Algorithm}
\makeatother

% ============================================================
% NOTATION — NUMBER THEORY
% ============================================================

\DeclareMathOperator{\lcm}{lcm}
\DeclareMathOperator{\ord}{ord}
\DeclareMathOperator{\Adv}{Adv}
\DeclareMathOperator{\poly}{poly}
\DeclareMathOperator{\negl}{negl}
\DeclareMathOperator{\Supp}{Supp}

% Sets and structures
\newcommand{\Zset}{\mathbb{Z}}
\newcommand{\Nset}{\mathbb{N}}
\newcommand{\Rset}{\mathbb{R}}
\newcommand{\Fset}{\mathbb{F}}
\newcommand{\Zn}[1]{\mathbb{Z}_{#1}}          % \Zn{n}  -> Z_n
\newcommand{\Zns}[1]{\mathbb{Z}_{#1}^{*}}     % \Zns{p} -> Z_p^*
\newcommand{\Fq}[1]{\mathbb{F}_{#1}}          % \Fq{q}  -> F_q
\newcommand{\bits}{\{0,1\}}
\newcommand{\bitstr}[1]{\{0,1\}^{#1}}         % \bitstr{n}

% ============================================================
% NOTATION — SCHEMES AND OPERATIONS
% ============================================================

\newcommand{\Gen}{\mathsf{Gen}}
\newcommand{\Enc}{\mathsf{Enc}}
\newcommand{\Dec}{\mathsf{Dec}}
\newcommand{\Sign}{\mathsf{Sign}}
\newcommand{\Vrfy}{\mathsf{Vrfy}}
\newcommand{\Mac}{\mathsf{Mac}}
\newcommand{\PRF}{\mathsf{F}}
\newcommand{\Hash}{\mathsf{H}}
\newcommand{\KDF}{\mathsf{KDF}}

\newcommand{\xor}{\oplus}
\newcommand{\concat}{\mathbin\Vert}
\newcommand{\getsr}{\stackrel{\scriptscriptstyle\$}{\leftarrow}} % x \getsr S

\newcommand{\secpar}{\lambda}
\newcommand{\adv}{\mathcal{A}}
\newcommand{\sect}[1]{\mathcal{#1}}
\newcommand{\msgspace}{\mathcal{M}}
\newcommand{\keyspace}{\mathcal{K}}
\newcommand{\ctspace}{\mathcal{C}}

\newcommand{\abs}[1]{\left\lvert #1 \right\rvert}
\newcommand{\ceil}[1]{\left\lceil #1 \right\rceil}
\newcommand{\floor}[1]{\left\lfloor #1 \right\rfloor}
\newcommand{\Prob}[1]{\Pr\!\left[ #1 \right]}

% ============================================================
% TIKZ — TWO-PARTY PROTOCOL DIAGRAMS
% ============================================================

\tikzset{
    party/.style        = {font=\bfseries},
    lifeline/.style     = {thick, rulegray, densely dashed},
    msg/.style          = {-{Stealth[length=2.6mm]}, thick, schoolred},
    msglabel/.style     = {font=\small, inner sep=2pt, fill=white},
    compute/.style      = {draw=rulegray, fill=black!3, rounded corners=2pt,
                           font=\small, inner sep=5pt, align=center}
}

% \begin{protocol}{caption}{label} ... tikz code ... \end{protocol}
\newenvironment{protocol}[2]
  {\def\protocolcaption{#1}\def\protocollabel{#2}%
   \begin{figure}[H]\centering\begin{tikzpicture}}
  {\end{tikzpicture}\caption{\protocolcaption}\label{\protocollabel}\end{figure}}

